The comparison reduces the problem to tracking repairs throughout the process. While a fixed fraction of the pieces remains, the
degrees and normalizers concentrate around their deterministic trajectories (Section~\ref{sec:regularidad}). A \emph{pattern} specifies
conditions on the squares from which queens are removed in a repair or its reverse and on the diagonals of the squares to be occupied.
Each specified square must already contain its queen or still be available, and each specified diagonal must be free of all pieces or
at least free of queens. For the forward count we track $32$ patterns for each candidate square of a crossing, and for the reverse
count we track $16$ for each pair of squares. Within each family, the expected change of a pattern can be expressed in terms of
patterns in the same family, up to an error bounded in the proof. Solving the corresponding linear system backward gives a linear
combination of patterns, the \emph{compensator}, that agrees with the counter at the chosen time and has small drift. A martingale
inequality then yields concentration up to an early cut-off, when about $\delta n$ queens remain to be placed. After the cut-off, each
queen changes the forward counter at each crossing by at most five and the reverse counter by $O(n)$. After division by $n$ and
$n^2$, respectively, the effect of the entire final stage is $O(\delta)$.

Once the loss after the cut-off has been bounded, the binary entropy inequality transfers concentration to the uniform law on each
layer. For the lower bound, the forward count suffices, since with probability tending to one every empty crossing admits at least
$a_0n/2$ repairs. To find the limit, we must take the position of the crossing into account. At $(r,c)$, the forward count concentrates
around $n\,a(\bar r,\bar c)$. Choosing the weight $w=1/a$ gives $D_{1/a}/(h^2n)\to1$. With this same weight, the mean of the reverse
counter divided by $n^2$ tends to $\int\Gamma/a$, where $\Gamma(r,c)$ is the limiting density of reverse operations that create the
crossing $(r,c)$ (Lemma~\ref{lem:nucleo}). The identity $\Gamma=b_0\,r_Rr_C\,a$, which follows from the product form of the profile,
gives $\int\Gamma/a=b_0$. Hence $f_h\to1/b_0=F(0)$. \ifverificacionLean Section~\ref{sec:verificacion} describes the Lean formalization
of the main statements and specifies what it verifies. \fi

\begin{figure}[t] \centering \begin{tikzpicture}[
 caja/.style={draw,rounded corners=2pt,align=center,font=\footnotesize,inner sep=3pt,
 minimum height=9mm,text width=29mm},
 flecha/.style={-{Stealth[length=2mm]},thick}]
\node[caja] (perfil) at (0,0) {Variational profile\\ and first-order bound (§\ref{sec:preliminares})};
\node[caja] (ref) at (3.6,0) {Entropy\\ comparison (§\ref{sec:referencia})};
\node[caja] (din) at (7.2,0) {Dynamic\\ regularity (§\ref{sec:regularidad})};
\node[caja] (abs) at (3.6,-1.7) {32 forward patterns\\ and early cut-off (§\ref{sec:absorbedores})};
\node[caja] (inv) at (7.2,-1.7) {16 reverse\\ patterns (§\ref{sec:inversa})};
\node[caja] (dc) at (0,-3.4) {Double counting\\ between layers (§\ref{sec:doblecuenta})};
\node[caja] (pos) at (3.6,-3.4) {Theorem~\ref{thm:positiva}\\ (§\ref{sec:positiva})};
\node[caja] (lim) at (7.2,-3.4) {Theorems~\ref{thm:limite} and~\ref{thm:estructura}\\ (§\ref{sec:limite})};
\node[caja] (cor) at (10.8,-3.4) {Corollaries~\ref{cor:cociente}\\ and~\ref{cor:capas} (§\ref{sec:corolarios})};
\draw[flecha] (perfil) -- (ref);
\draw[flecha] (ref) -- (din);
\draw[flecha] (din) -- (abs);
\draw[flecha] (din) -- (inv);
\draw[flecha] (abs) -- (pos);
\draw[flecha] (abs) -- (lim);
\draw[flecha] (inv) -- (lim);
\draw[flecha] (dc) -- (pos);
\draw[flecha] (dc.south) to[out=-25,in=-155] (lim.south);
\draw[flecha] (lim) -- (cor);
\end{tikzpicture} \caption{Logical structure of the proof. The double count is purely combinatorial. Every arrow leading to a result
under $\mu_h$ uses the entropy comparison of Section~\ref{sec:referencia}.} \label{fig:esquema} \end{figure}

\subsection{Related work}

Simkin~\cite{Simkin} characterizes the exponential term of $Q(n)$ through entropy maximization and proves a large deviation principle
for the profile of complete solutions. Earlier, Luria~\cite{Luria} obtained upper bounds using the entropy method. The first-order
lower bound we use comes from Simkin's result (Theorem~\ref{thm:simkin}). To prove it, Simkin fixes a step \emph{queenon}, that is, a
queenon (Section~\ref{sec:preliminares}) whose density is constant on each block of a grid. At each step he chooses a block with the
probability assigned to it by this measure and places a queen at random on one of its available squares. At the end he
completes the placement using absorbers~\cite[Section~5]{Simkin}. Luria and Simkin~\cite{LS} use a random greedy process to construct
near-solutions, which they complete using absorbers to obtain a counting lower bound. Bowtell and Keevash~\cite{BK} study enumeration
on the torus and its consequences for the ordinary board. The reference process and the absorbers of this paper are in the same spirit, but
success under a construction process does not by itself identify the uniform law on placements with $h$ holes; that identification
is the role of the entropy comparison. Comparing the uniform law with the output law of a random process is the main tool in Kwan's
work~\cite{Kwan} on Steiner triple systems. Double counting the repairs between consecutive layers is an instance of the method of
\emph{switchings}~\cite{McKayWormald,HMcK}.

The problem of completing a given partial placement dates back to 1850 and is NP-complete~\cite{GJN}. Glock, Munhá Correia, and
Sudakov~\cite{GMS} prove that every legal placement of at most $n/60$ queens can be completed and construct placements of size close to
$n/4$ that admit no completion. Nielsen~\cite{Nielsen} improves the upper bound on the extremal threshold to $0.216\,n$ for
sufficiently large $n$. These are worst-case results. Here we study probabilities under the uniform law on nearly complete placements.

To the best of our knowledge (our literature search was not exhaustive), Theorems~\ref{thm:limite} and~\ref{thm:estructura} and
Corollary~\ref{cor:cociente} are new. We make no claim of novelty for the general tools used here, namely the chain rule for relative
entropy, comparison with random processes~\cite{Kwan}, double counting and switchings~\cite{HMcK}, the differential equation
method~\cite{Wormald}, and concentration inequalities~\cite{AlonSpencer}.
